% Einheit 2 von 5 – Eigenschaften aller Graphen am Term (bank/funktionsscharen-und-ortskurven/e2.jsonl, zusammenbau v0.1)
%% TODO Zeitmarke Einheit 2: Marken-Zeile nicht lesbar
%% TODO Zweigzeile Teil 1: was hier gelernt wird (Satz in Schülersprache aus dem Einheitstitel) – Einheit 2
\einheitenkopf[e2]{Einheit 2 von 5 · Eigenschaften aller Graphen am Term}
\zweigzeile{Abitur LK}
\setcounter{aufgabe}{11}

%% TODO Titel: Ich-kann-Satz fehlt in der Bank (Platzhalter: Kettenname) – Nr. 12
%% TODO Anweisung: ein Satz über der ersten Teilaufgabe fehlt in der Bank – Nr. 12
\begin{aufgabe}{Eigenschaften aller Graphen}
%% TODO \rechenplatz{n}: Schrittzahl des Grundfalls fehlt in der Bank (nur bei fester Schreibform, 2.3 b)
\begin{teile}
\teil Setze $x = 2$ in $f_a(x) = a \cdot (x - 2) + 5$ ein. Fällt der Parameter $a$ heraus? \\ \kreuz{fällt heraus} \kreuz{bleibt}
\teil Gib die Nullstellen von $f_a(x) = a x \cdot (x - 6)$ mit $a \neq 0$ an. $x_1 =$ \_\_, $x_2 =$ \_\_
\teil Gegeben ist $f_a(x) = x^3 \cdot e^{a x}$ mit $a \neq 0$. Begründe, dass der Graph für $x < 0$ unterhalb der x-Achse verläuft. \hfill (Abitur 2024 LK)
\teil Gegeben ist $f_a(x) = (x^2 - a) \cdot e^{2 - x}$ mit $a \in \mathbb{R}$. Gib die Anzahl der Nullstellen in Abhängigkeit von $a$ an. \hfill (Abitur 2018 LK)
\teil Gegeben ist $f_a(x) = (2x + 1) \cdot e^{a x}$ mit $a \in \mathbb{R}$. Wie verhalten sich die Funktionswerte für $x \to +\infty$? Unterscheide nach $a$. \hfill (Abitur 2024 LK)
\teil Gegeben ist $f_r(x) = \frac{2}{r} \cdot (x^2 - 5x) + 3$ mit $r \neq 0$. Bestimme die Punkte, durch die alle Graphen verlaufen. \hfill (Abitur 2018 LK)
\teil Gegeben ist $f_a(x) = x^3 \cdot e^{a x^2}$ mit $a \neq 0$. Weise nach, dass jeder Graph der Schar punktsymmetrisch zum Ursprung ist. \hfill (Abitur 2026 LK)
\teil Gegeben ist $f_a(x) = 2x \cdot e^{a x}$ mit $a \in \mathbb{R}$. Weise nach, dass alle Graphen im Ursprung dieselbe Steigung haben. \hfill (Abitur 2025 LK)
\teil Für die Schar $g_k$ gilt für alle $k$, $k_1$ und $k_2$: (1) $g_k(1) = 2$; (2) $g_k'(1) = g_0'(1)$; (3) $g_{k_1}(x) = g_{k_2}(x) \Leftrightarrow k_1 = k_2$ oder $x = 1$. Was folgt daraus für den Verlauf der Graphen? \hfill (Abitur 2022 LK)
\end{teile}
\end{aufgabe}

%% TODO Titel: Ich-kann-Satz fehlt in der Bank (Platzhalter: Kettenname) – Nr. 13
%% TODO Anweisung: ein Satz über der ersten Teilaufgabe fehlt in der Bank – Nr. 13
\begin{aufgabe}{Eigenschaften aller Graphen – Fehler finden, Begründen}
\begin{teile}
\teil Lea bestimmt die Nullstellen von $f_a(x) = (x - a) \cdot e^{x + 1}$: „$x = a$ oder $e^{x + 1} = 0$, also $x = -1$.“ Finde den Fehler und gib die Nullstellen richtig an.
\teil Begründe, warum an einem gemeinsamen Punkt aller Graphen der Schar $f_a(x) = a x^2 + x$ der Parameter herausfallen muss.
\end{teile}
\end{aufgabe}
